\section{Discussion}
\label{sec:discussion}

The experimental evaluation of the \textit{MVC Test Orchestrator}
across 20 distinct SRS documents indicates that the tested pipeline can support selected MVC conformance tasks in newly generated Python artifacts under the studied configuration. The following
subsections discuss these findings within the broader context of
software maintainability and architectural conformance,
organized around four axes that correspond to the five Research
Questions: semantic coverage (RQ1), rule-based control (RQ2 \& RQ3), 
human-AI agreement (RQ4), and stochastic stability (RQ5).

\subsection{Requirement Synthesis and Semantic Coverage (RQ1)}

The initial extraction phase represents the system's capacity to translate unstructured textual inputs into a structured, computable architectural structure. While the binary Requirement Coverage metric utilized in our preliminary framework~\citep{dim2026mvc} was insufficient for capturing complex linguistic variances, the integration of BERTScore in the present study provides quantitative support for the system's semantic fidelity. 

The consistently high BERTScore values demonstrate that the LLM does not merely duplicate functional specifications from the SRS document; rather, it successfully synthesizes these definitions into the \texttt{requirements\_analysis.json} artifact (Table~\ref{tab:artifact_flow}) while faithfully preserving domain-specific technical terminology. This indicates that the LLM's generative paraphrasing capabilities are more likely to reflect semantic synthesis than stochastic hallucination under the conditions tested in this study. Ultimately, these findings suggest that transformer-based semantic verification can help reduce the risk of "information loss" during the critical transition from requirements engineering to architectural design.


\subsection{Rule-Based Control in Combating ``Vibe Coding'' (RQ2 \& RQ3)}

Current literature indicates that unconstrained LLMs systematically violate macro-level layered architectural paradigms, often introducing structural flaws at notably high rates. The \textit{MVC Test Orchestrator} distinguishes itself from baseline generative tools through its Architecture-First governance model (Section~\ref{sec:introduction}).

\paragraph{Contextualizing Detection Precision (RQ2):}
Prior defect-detection studies address different tasks, datasets, labels, or success oracles than the fault-injection benchmark used in this study. Their reported percentages are therefore not directly comparable with the ACP and VDR figures reported here; we cite these works only to position the problem and evaluation constructs, and no performance superiority is inferred from cross-study numerical differences.

The AST-based static analyzer demonstrated consistent reliability, achieving 96.2\% precision and 98.8\% recall (F1 = 97.5\%) on the fault-injection benchmark. Recent state-of-the-art approaches, such as the DeepTriFix framework proposed by \citet{wan2025deeptrifix}, adopt a hybrid architecture combining ASTs with LLMs for code defect detection; when applied to generalized vulnerability datasets, such systems report detection accuracies around 67.6\%~\citep{wan2025deeptrifix}. The difference in scope explains the difference in reported figures: frameworks like DeepTriFix are engineered to identify complex, open-ended operational vulnerabilities (e.g., memory leaks) across highly variable semantic contexts, whereas our module functions as a dedicated architectural auditor that measures source code solely against the structured, closed-loop constraints of the \textit{Architectural Map}. By narrowing the detection scope from universal logic errors to specific MVC structural violations, the system reduces the false-positive surface that typically affects generalized semantic analyzers.

\paragraph{Contextualizing Remediation Success (RQ3):}
Prior program-repair studies address different tasks, datasets, and success oracles than the architectural re-audit criterion used in this study. Their reported percentages are therefore not directly comparable with the RSDA and RSR\textsubscript{e2e} figures reported here; we cite these works only to position the problem and evaluation constructs, and no performance superiority is inferred from cross-study numerical differences.

Following detection, the autonomous self-healing capability yielded 92.6\% repair success among detected violations (RSDA), and 91.6\% end-to-end (RSR\textsubscript{e2e}). To contextualize the design rationale, a recent evaluation of a standalone code-specialized LLM (DeepSeek-Coder-33B) reports repair rates between 5.7\% and 21.9\% on enterprise-grade bugs, and around 45.5\% in simplified algorithmic scenarios~\citep{sun2025empirical}. These figures reflect a different repair scope and oracle, and are cited here to illustrate why a constrained, architecture-aware approach was chosen over unconstrained generation. Instead of relying on open-ended LLM generation, the \textit{Fixer Agent} operates under defined architectural boundaries, channeling the LLM's reasoning toward localized logic patching and reducing the search space.

\subsection{Methodological Validity of Human-AI Agreement (RQ4)}
\label{sec:discussion_rq4}

To comprehensively address \textbf{RQ4}, the autonomously generated architectural designs were subjected to a heuristic evaluation via the Architectural Component Extraction Score (ACES) metric. This validation was conducted by a panel of five expert software engineers (N=5) utilizing a structured 5-point Likert scale. In empirical software engineering, expert-based evaluation by a focused panel of domain experts is a well-established methodology for assessing system prototypes and outputs~\citep{monasor2014walk}.

As demonstrated by \citet{monasor2014walk}, a well-designed heuristic evaluation relies on a small, targeted number of expert participants to provide an assessment that is both deep and broad, particularly when utilizing structured Likert scales to measure specific qualitative and functional dimensions. By adopting this empirical standard, our framework avoids the superficiality of mass-crowdsourced evaluations. Furthermore, to mitigate the subjective bias inherent in human evaluation, an Inter-Rater Reliability (IRR) analysis was conducted. Inter-rater reliability was assessed using a two-way random-effects, absolute-agreement, average-measures intraclass correlation coefficient, ICC(A,5), because the reported score is the mean of five raters and the panel is treated as a sample from a broader population of qualified experts, computed at the criterion level (SRS $\times$ MVC layer $\times$ rubric criterion, 240 rating targets). The estimate was ICC = 0.942, 95\% CI [0.93, 0.95], indicating excellent agreement among the five expert raters and supporting the reliability of the reported ACES scores. The consensus achieved by the N=5 expert panel, resulting in an ACES score of 74.0\%, is not merely a statistical average, but a strong, methodologically sound indication that the generated MVC architectures align closely with professional software engineering expectations. The layer-specific score distributions in Table~\ref{tab:rq1_results} further qualify this result: Model adequacy was the lowest and most variable layer (mean 2.82/5, SD\,$=$\,0.51), while View adequacy was consistently higher (mean 4.41/5, SD\,$=$\,0.37). Causal attribution is not possible from the current data alone; no requirement-level error taxonomy was collected.

\subsection{Stochastic Mitigation and Architectural Stability (RQ5)}

A fundamental challenge in employing Large Language Models for software engineering is their built-in stochastic variability. To address \textbf{RQ5}, our multi-iteration analysis yielded an overall Consistency Score (CS) of 70.6\%. Prior structured-output consistency studies address different tasks, output formats, and measurement constructs than the CS metric used in this study. Their reported percentages are therefore not directly comparable with the CS figures reported here; we cite this work only to position the problem, and no performance superiority is inferred from cross-study numerical differences.

Recent benchmarking by \cite{wang2025sted} on structured output reliability reports that standalone LLMs (such as Llama-3.3-70B and various Claude models) can exhibit consistency scores between 41.3\% and 65.8\% under standard operational temperatures. These figures reflect a different output domain and consistency construct, and are cited here to illustrate the general challenge of LLM stochasticity that motivated our rule-based approach. Our framework's 70.6\% overall CS, with core Model extractions regularly achieving near-perfect stability, indicates that the LLM-driven requirement-extraction and architectural-design phases (Phases I–II) exhibited moderate-to-high run-to-run stability under the tested Gemini 2.5 Flash configuration. Because CS was measured only on these LLM-generated artifacts, this result reflects the underlying model's own stability rather than an ablation-verified effect of the downstream rule-based stages (Scaffolder, AST auditor); generalization beyond the recorded configuration requires further investigation.
